% Figure: the Rayleigh residue curve for a simple batch (differential)
% distillation at constant relative volatility. The still-pot liquid mole
% fraction of the light component falls as the charge is boiled away, and the
% curve plots the fraction of the original charge still in the pot against the
% still composition. A more volatile light component (larger alpha) empties
% the light component from the pot faster, so its curve drops away from the
% starting composition sooner. The worked example reads its still composition
% off the alpha = 2.5 curve at a boil-off of one third of the charge.
\begin{figure}[htbp]
\centering
\begin{tikzpicture}
\begin{axis}[
  width=11cm, height=6.8cm,
  xlabel={fraction of the original charge remaining in the still $W/W_0$},
  ylabel={still-pot light-component mole fraction $x$},
  xmin=0, xmax=1, ymin=0, ymax=0.8,
  legend pos=south east,
  legend cell align=left,
  grid=both, grid style={enggray!18},
  samples=120, domain=0.02:1,
]
% Rayleigh integral for constant alpha, started at x0 = 0.5.
% ln(W/W0) = integral_{x0}^{x} dx / (y - x), y = alpha x /(1+(alpha-1)x).
% For constant alpha this integrates in closed form. We plot x as the
% parameter and W/W0 = exp(F(x)-F(x0)) with
% F(x) = (1/(alpha-1)) [ ln(x/(1-x)) - alpha ln x ... ] handled numerically
% via the standard closed form:
% W/W0 = [ (x/x0) ]^{1/(alpha-1)} * [ (1-x0)/(1-x) ]^{alpha/(alpha-1)} .
% alpha = 2.5, x0 = 0.5
\addplot[engblue, very thick, domain=0.02:0.5]
  ( { (x/0.5)^(1/1.5) * (0.5/(1-x))^(2.5/1.5) }, {x} );
\addlegendentry{$\alpha = 2.5$, $x_0 = 0.5$}
% alpha = 5, x0 = 0.5 (more volatile light component empties faster)
\addplot[engred, very thick, dashed, domain=0.02:0.5]
  ( { (x/0.5)^(1/4) * (0.5/(1-x))^(5/4) }, {x} );
\addlegendentry{$\alpha = 5$, $x_0 = 0.5$}
% starting point marker
\addplot[engblack, only marks, mark=*, mark size=1.6pt] coordinates {(1,0.5)};
\node[font=\scriptsize, engblack, anchor=south east] at (axis cs:0.98,0.52)
  {start of the boil};
% worked-example operating point on the alpha=2.5 curve at W/W0 = 2/3
\addplot[engorange, only marks, mark=square*, mark size=1.8pt]
  coordinates {(0.667,0.408)};
\node[font=\scriptsize, engorange, anchor=west] at (axis cs:0.62,0.30)
  {worked example};
\end{axis}
\end{tikzpicture}
\caption[Rayleigh residue curve for a batch distillation]{The still-pot
composition falling as a simple batch distillation boils the charge away, at
two relative volatilities. The pot starts full at the right, at the charge
composition, and empties leftward as vapour is drawn off. A more volatile
light component, the larger relative volatility, strips the light component
out of the pot faster, so its curve peels down from the starting composition
sooner. The worked example reads the still composition off the lower-volatility
curve after one third of the charge has boiled over, the point marked. A batch
still has no reflux and no operating line, so the Rayleigh integral, not a
McCabe-Thiele staircase, sets its performance.}
\label{fig:vol04ch07:rayleigh-batch}
\end{figure}
